\section{Open Research Problems}

\label{reserch_problems}






Aside from its utility for model evaluations, \our{} also offers a test-bed for rapidly prototyping new inference-time methods for improving the agentic capabilities of LLMs and VLMs. There are many open research problems in this space. As of the writing of this paper, some of the most performant methods for improving model reasoning capabilities on short-form and/or shorter-context problems are infeasible to apply naively to \our{} due to the extremely long-context nature of tasks. Addressing these challenges could further enhance the development of stronger autonomous agents. We highlight several key areas for future work below.

\paragraph{In-Context Learning and Few-Shot Prompting} \our{} enables evaluation of In-Context Learning (ICL) agents, which can use few-shot examples to adapt to out-of-distribution tasks. We provide a small dataset of human demonstrations for each environment and an implementation of few-shot conditioning in the \our{} codebase. The benchmark codebase also supports the study of In-Context Reinforcement Learning \citep{lee2024supervised, laskin2022context, lin2023transformers}, where agents learn to improve from mistakes during inference. On the large models benchmarked in Section \ref{section:results}, naive few-shot learning (i.e., prompting LLM and VLM agents with examples of full human games in-context) is extremely computationally expensive to run on \our{}. For example, a single demonstration of NetHack game-play can require upwards of $700,000$ input tokens to represent in a prompt. Despite advancements in fast inference technologies like caching and falling API costs for long-context prompting, we found these experiments to be infeasible to conduct at this time. Sub-selecting only the relevant parts of demonstrations via retrieval-augmented few-shot prompting strategies~\citep{lewis2020retrieval} might offer a way to circumvent these challenges. We leave exploration of such methods for future work.

\paragraph{Advanced Reasoning Strategies} Beyond simply prompting LLMs and VLMs to directly predict the next action of game-play, \our{} also supports the study of more advanced reasoning techniques like chain-of-thought \citep{wei2022chain}, self-refinement \citep{madaan2024self}, and basic planning. These methods have been demonstrated to improve model performance on shorter-context problems. We believe them to be an exciting direction for future work on long-context reasoning and decision-making. For example, model performance on the tasks in \our{} might be improved by integrating multi-agent collaboration \citep{chang2023prompting, khan2024debating, yao2024tree} and tool usage \citep{shen2024hugginggpt, ruan2023tptu, schick2024toolformer, qin2023toolllm} in decision-making. Additionally, incorporating memory mechanisms or reinforcement learning techniques could help bridge the "knowing-doing" gap, enabling models to apply their knowledge effectively in practical, long-horizon tasks. Finally, experimenting with open-ended self-improvement loops \citep{wang2023voyager, hu2024automated} could lead to more adaptive and general agents \citep{team2023human, hughes2024open}, offering a pathway toward truly autonomous systems.

\paragraph{Limitations of Current Vision-Language Models} Despite their potential, our benchmark shows significant variability in VLM performance. While some models, like Llama 3.2, struggle to integrate visual information into coherent decision-making, others-most notably Sonnet 3.5-demonstrate stronger performance in VLM mode. This disparity highlights significant variability in VLM capabilities, which may stem from differences in training objectives and datasets. For example, Sonnet 3.5's superior performance can be attributed in part to its training on tasks involving computer usage \citep{anthropic2024developing}, which inherently require integrating visual and textual inputs for action-based reasoning.

Recent studies have identified key limitations of VLMs that align with our findings, including biases toward natural image-text pairs, optimization for image description rather than action-oriented reasoning, and challenges with out-of-distribution inputs~\citep{tan2024towards, tong2024eyes, rahmanzadehgervi2024vision, zang2024overcoming, guan2023hallusionbench}. These limitations are further exemplified in our benchmark, where grid-based image observations differ significantly from the natural image-text pairs on which many VLMs are trained~\citep{yu2023scaling, rahmanzadehgervi2024vision}. Moreover, the computational cost of image processing constrained our evaluation to a single image per observation, with the remainder of the history provided in text. While this constraint may hinder performance for some models, our results show that certain VLMs like Claude 3.5 Sonnet can still perform robustly under these conditions.

To address these challenges, our codebase already supports multi-image observation histories, and future iterations will incorporate video observations, which are likely better suited for the long-horizon sequential decision-making tasks central to our benchmark. These enhancements aim to better evaluate and leverage the potential of VLMs in complex reasoning scenarios. We plan to introduce support for video observations once prominent models with efficient video-processing capabilities become available, ensuring that our benchmark remains aligned with the latest advancements in VLM technology.




\paragraph{Computational Limitations of Large Language Models} Mechanistic interpretability could provide valuable insights for understanding the computational limitations of agentic LLMs. The computational expressiveness of LLMs is fundamentally linked with the ability to solve complex reasoning problems ~\citep{wei2022emergent}.
While current models perform well on simple tasks such as navigation and object manipulation, they struggle with more complex tasks that could require non-trivial and general-purpose computation, for example, building a shelter or developing combat strategies. 
This could be due to the models' inability to retrieve relevant computational circuits~\citep{olah2020zoom},  limitations to inference-time budget~\citep{snell2024scaling}, or representational expressivity. 
This raises important questions about the scope of effectively solvable tasks for LLMs and VLMs, which is dependent on factors such as model depth, context size, and the distribution shift between pre-training and downstream tasks. Further research is needed to understand the underlying causes of these limitations and to develop strategies for overcoming them, such as adaptive simulation of computational circuits during runtime. 








